% This document was generated by an LLM.
% Content fragment; compile intro-analysis-self-check.tex.
% Coverage: Hunter (2012), Chapter 3, printed pp. 21--37.
% Corrections and clarifications:
% - Sequential continuity also holds at isolated domain points (Theorem 3.6).
% - In Example 3.14, the rational branch excludes 0; f(0)=0 by convention.
% - In the cover proof of Theorem 3.31, preimages are relatively open in K;
%   lift them to ambient open sets before applying Chapter 1's definition.
% - Theorem 3.33 requires a nonempty domain to assert attained extrema.
% - In Example 3.37, the maximum and upper bound are 4/pi, not 1/pi.
% - The endpoint-image equality in Example 3.46 can hold for some functions
%   that are not increasing; increasingness is sufficient, not necessary.

\section{Continuous Functions}

\subsection{Continuity}
\textit{Source: \S3.1, pp.~21--25.}

\begin{enumerate}
  % Definition 3.1.
  \item State the \(\epsilon\)-\(\delta\) definition of continuity
    of \(f:A\to\R\) at \(c\). Define continuity on \(B\subseteq A\)
    and continuity of a function on its domain.

  % Definition 3.2 and the exercise following it.
  \item State the neighborhood definition of continuity at a point.
    Explain its equivalence with the \(\epsilon\)-\(\delta\) definition.
    How does it differ from the neighborhood definition of a limit?

  % Discussion of isolated points following Definition 3.2.
  \item Prove that every function is continuous at each isolated point
    of its domain. Why is no accumulation-point hypothesis needed
    in the definition of continuity?

  % Discussion at the start of p. 22 and Example 3.3.
  \item At an accumulation point \(c\in A\), express continuity in
    terms of a function limit. What does this imply for a function
    defined on an open interval?

  % Example 3.4.
  \item Characterize continuity of \(f:[a,b]\to\R\), with \(a<b\),
    using limits at interior points and one-sided limits at the endpoints.

  % Example 3.5.
  \item Let \(A=\{0\}\cup\{1/n:n\in\N\}\), and define
    \(f(0)=y_0\), \(f(1/n)=y_n\). At which points is continuity
    automatic? Give a necessary and sufficient condition on \((y_n)\)
    and \(y_0\) for \(f\) to be continuous on \(A\).

  % Theorem 3.6, extended to isolated domain points.
  \item State the sequential characterization of continuity
    (Theorem 3.6). May the testing sequences have terms equal to \(c\)?
    Explain why the characterization also holds if \(c\) is isolated.

  % Observation following Theorem 3.6.
  \item How can a single sequence establish discontinuity at a point?
    State the conditions that the sequence and its function values
    must satisfy.

  % Example 3.7.
  \item Prove directly from the definition that \(f(x)=\sqrt{x}\)
    is continuous on \([0,\infty)\). Give explicit choices of
    \(\delta\), treating \(c>0\) and \(c=0\) separately.

  % Example 3.8.
  \item Use a trigonometric identity and \(|\sin t|\leq|t|\) to prove
    that sine is continuous at every real point. Can the same choice
    of \(\delta\) work at every point for a fixed \(\epsilon\)?

  % Example 3.9.
  \item Determine all continuity points of \(\operatorname{sgn}x\).
    What kind of discontinuity does it have at zero, and can changing
    its value at zero make it continuous there?

  % Example 3.10.
  \item Define \(f(0)=0\) and \(f(x)=1/x\) for \(x\ne0\).
    Determine its continuity points and describe its one-sided behavior
    at zero. Distinguish divergence to infinity from a finite one-sided limit.

  % Example 3.11.
  \item Define \(f(0)=0\) and \(f(x)=\sin(1/x)\) for \(x\ne0\).
    Determine its continuity points and show why no choice of \(f(0)\)
    can make it continuous at zero.

  % Example 3.12 and its following discussion.
  \item Define \(f(0)=a\) and \(f(x)=x\sin(1/x)\) for \(x\ne0\).
    For which \(a\in\R\) is \(f\) continuous at zero?
    Prove your answer and classify the discontinuity for the other values.

  % Example 3.13.
  \item Define the Dirichlet function. Prove that it is discontinuous
    everywhere and has no limit at any real point. Where do density
    of the rationals and density of the irrationals enter the argument?

  % Example 3.14; specify the modified value at zero unambiguously.
  \item Define the version of the Thomae function used in the notes,
    with value zero at zero and at irrational points, and value \(1/q\)
    at a nonzero rational \(p/q\) in lowest terms with \(q>0\).
    State exactly where it is continuous and where it is discontinuous.

  % Classification at the end of Section 3.1.
  \item Define removable, jump, and essential discontinuities using
    finite limits, following the terminology in these notes.
    Give an example of each. Which can be repaired by changing the
    function's value at the discontinuity?
\end{enumerate}

\subsection{Properties of continuous functions}
\textit{Source: \S3.2, pp.~25--26.}

\begin{enumerate}[resume]
  % Theorem 3.15.
  \item State the continuity rules for scalar multiples, sums,
    products, and quotients (Theorem 3.15). Include the extra condition
    for a quotient and explain how the rules follow from limit laws
    at accumulation points. What happens at isolated domain points?

  % Definitions preceding Corollary 3.16.
  \item Define polynomial and rational functions with real coefficients.
    What is the natural domain of a rational function given by \(P/Q\)?

  % Corollary 3.16 and its proof.
  \item State and prove the continuity result for polynomials and
    rational functions (Corollary 3.16). Which elementary continuous
    functions and algebraic rules suffice?

  % Example 3.17.
  \item Determine the domain and continuity points of
    \[
      f(x)=\frac{x+3x^3+5x^5}{1+x^2+x^4}.
    \]
    Justify your answer without an \(\epsilon\)-\(\delta\) computation.

  % Definition of composition preceding Theorem 3.18.
  \item Define \(g\circ f\) for \(f:A\to\R\) and \(g:B\to\R\).
    What relation between \(f(A)\) and \(B\) makes the composition
    well-defined on all of \(A\)?

  % Theorem 3.18 and Corollary 3.19.
  \item State the pointwise and whole-domain continuity theorems for
    compositions (Theorem 3.18 and Corollary 3.19). At which point
    must the outer function be continuous when testing continuity at \(c\)?

  % Proof of Theorem 3.18.
  \item Prove continuity of a composition from the definition.
    Explain the order in which you choose the output tolerance,
    the intermediate tolerance, and the input radius.

  % Examples 3.20--3.21.
  \item Use composition and product rules to justify continuity of
    \(\sin(1/x)\) and \(x\sin(1/x)\) at every nonzero point.
    Why do these rules alone say nothing about continuity at zero?
\end{enumerate}

\subsection{Uniform continuity}
\textit{Source: \S3.3, pp.~27--29.}

\begin{enumerate}[resume]
  % Definition 3.22.
  \item Define uniform continuity on \(A\). Compare the quantifier
    order with continuity at every point of \(A\). Which dependencies
    of \(\delta\) are permitted in each definition?

  % Discussion following Definition 3.22.
  \item Prove that uniform continuity implies continuity at every
    domain point. Is the converse true?

  % Proposition 3.23.
  \item Negate the definition of uniform continuity. State the
    equivalent sequential criterion for failure of uniform continuity
    (Proposition 3.23), including the fixed positive separation in
    function values. Must the two input sequences converge individually?

  % Proof of Proposition 3.23.
  \item Prove Proposition 3.23 in both directions. How does choosing
    \(\delta=1/n\) produce the required sequences?

  % Example 3.24.
  \item Prove that sine is uniformly continuous on \(\R\) by giving
    a single \(\delta\) for each \(\epsilon>0\).

  % Example 3.25.
  \item Prove that \(x\mapsto x^2\) is uniformly continuous on
    \([0,1]\), and then on every bounded subset of \(\R\).
    Give an explicit uniform choice of \(\delta\).

  % Example 3.26.
  \item Show that \(x\mapsto x^2\) is not uniformly continuous on
    \(\R\) by constructing two input sequences whose distances tend
    to zero while their squared values remain separated.

  % Example 3.27.
  \item Give a continuous function on the bounded interval \((0,1]\)
    that is not uniformly continuous. Verify failure using two
    explicit sequences from that interval.

  % Example 3.28.
  \item Give a bounded continuous function on \((0,1]\) that is not
    uniformly continuous. Construct sequences that establish the
    failure. Why does boundedness of both the domain and the range
    not suffice?
\end{enumerate}

\subsection{Continuous functions and open sets}
\textit{Source: \S3.4, pp.~29--30.}

\begin{enumerate}[resume]
  % Definitions at the beginning of Section 3.4.
  \item Define the image \(f(B)\) and preimage \(f^{-1}(C)\) of sets
    under \(f:A\to\R\). Does defining a preimage require \(f\)
    to have an inverse function?

  % Example 3.29.
  \item For \(f(x)=x^2\) on \(\R\), compute the image and preimage
    of \((1,4)\), and the image and preimage of \((-1,1)\).

  % Discussion following Example 3.29.
  \item Must a continuous function map open sets to open sets?
    Give a counterexample and distinguish images from preimages.

  % Theorem 3.30.
  \item State the characterization of continuity by preimages of open
    sets (Theorem 3.30). In which space must each preimage be open?
    Does the characterization concern continuity at one point or
    throughout the domain?

  % Forward direction of Theorem 3.30.
  \item Prove that a continuous \(f:A\to\R\) has relatively open
    preimages of open subsets of \(\R\). Construct a relative
    neighborhood at each point of a preimage.

  % Converse direction of Theorem 3.30.
  \item Prove the converse by taking preimages of
    \(\epsilon\)-neighborhoods of \(f(c)\).
    How does relative openness provide the required input radius?
\end{enumerate}

\subsection{Continuous functions on compact sets}
\textit{Source: \S3.5, pp.~30--32.}

\begin{enumerate}[resume]
  % Theorem 3.31.
  \item State the theorem on continuous images of compact sets
    (Theorem 3.31). Does the image have to be an interval?

  % Sequential proof of Theorem 3.31.
  \item Prove Theorem 3.31 using sequences. Given a sequence in
    \(f(K)\), how do you obtain a sequence in \(K\), and where do
    compactness and continuity enter the subsequence argument?

  % Open-cover proof of Theorem 3.31, with relative openness clarified.
  \item Prove Theorem 3.31 using open covers and preimages.
    The preimages are open in \(K\); how can they be written as
    intersections with open subsets of \(\R\) so that the ambient
    open-cover definition of compactness applies?

  % Example 3.32.
  \item Give a continuous function on a closed subset of \(\R\)
    whose image is not closed. Compute the image and identify the
    missing hypothesis that prevents Theorem 3.31 from applying.

  % Theorem 3.33; nonemptiness made explicit.
  \item State the Weierstrass extreme value theorem (Theorem 3.33),
    including the nonemptiness assumption needed for attained extrema.
    Distinguish boundedness from attainment of a supremum and infimum.

  % Proof of Theorem 3.33.
  \item Deduce the extreme value theorem from compactness of the
    continuous image. Why do the extrema of the image correspond
    to values attained at points of the domain?

  % Example 3.34.
  \item Give an unbounded function on a compact interval. Does it
    have a minimum or a maximum? Explain which hypothesis of the
    extreme value theorem fails.

  % Example 3.35.
  \item Give a continuous, unbounded function on a bounded interval.
    Does it have a minimum or a maximum? Explain which hypothesis
    of the extreme value theorem fails.

  % Example 3.36.
  \item Give a bounded continuous function that attains neither its
    supremum nor its infimum. Specify the domain and compute both bounds.

  % Example 3.37, correcting its reported upper bound and maximum.
  \item Define \(f:[0,2/\pi]\to\R\) by \(f(0)=0\) and
    \(f(x)=x+x\sin(1/x)\) for \(x>0\).
    Show that it is continuous and find its maximum and minimum values.
    Exhibit infinitely many points where the minimum is attained.
    Check the printed maximum in Example 3.37 against the value at
    the right endpoint.

  % Theorem 3.38.
  \item State the uniform continuity theorem for continuous functions
    on compact sets (Theorem 3.38). Which failures of compactness
    occur in the counterexamples of \S3.3?

  % Proof of Theorem 3.38.
  \item Prove Theorem 3.38 by contradiction using the sequential
    criterion for failure of uniform continuity. Why do the two
    selected subsequences converge to the same point, and why does
    that point belong to the domain?

  % Example 3.39.
  \item Explain why the oscillatory function in Example 3.37 is
    uniformly continuous on \([0,2/\pi]\), despite oscillating
    infinitely often near zero. Is an explicit uniform radius needed
    to apply Theorem 3.38?
\end{enumerate}

\subsection{The intermediate value theorem}
\textit{Source: \S3.6, pp.~32--35.}

\begin{enumerate}[resume]
  % Theorem 3.40.
  \item State the zero-existence version of the intermediate value
    theorem (Theorem 3.40), including the endpoint sign conditions
    and the location of the zero.

  % Proof of Theorem 3.40.
  \item Prove Theorem 3.40 for \(f(a)<0<f(b)\) using
    \(E=\{x\in[a,b]:f(x)<0\}\) and \(c=\sup E\).
    Why does \(c\) exist? Show why both \(f(c)<0\) and \(f(c)>0\)
    are impossible, and why the resulting zero lies in \((a,b)\).

  % Example 3.41.
  \item Give a continuous function on a compact set with positive
    and negative values but no zero. Explain why the interval
    hypothesis in Theorem 3.40 matters.

  % Example 3.42.
  \item Give a function on a compact interval with opposite signs at
    the endpoints but no zero. Explain why continuity cannot be
    omitted from Theorem 3.40.

  % Example 3.43.
  \item Use the intermediate value theorem to prove existence of a
    positive real solution of \(x^2=2\). What additional argument
    gives uniqueness of the positive solution?

  % Bisection discussion following Example 3.43.
  \item Describe the bisection method for a continuous function with
    opposite endpoint signs. Apply the first two bisections to
    \(x^2-2\) on \([1,2]\). How does the enclosing interval's
    length change after \(n\) steps if no zero is found exactly?

  % Completeness discussion following Example 3.43.
  \item Restrict \(f(x)=x^2-2\) to \([1,2]\cap\Q\).
    Why is this function continuous on its domain but nowhere zero?
    What does the example show about the role of real completeness
    in the intermediate value theorem?

  % Theorem 3.44 and its proof.
  \item State the general intermediate value theorem (Theorem 3.44)
    and deduce it from the zero-existence version by shifting the
    function by a constant. How do you handle the two possible orders
    of the endpoint values?

  % Theorem 3.45.
  \item State the result describing the image of a compact interval
    under a continuous real-valued function (Theorem 3.45).
    How are the endpoints of the image determined?

  % Proof of Theorem 3.45.
  \item Prove Theorem 3.45 using the extreme value and intermediate
    value theorems. Explain separately why the image is contained
    in the proposed interval and why it contains that entire interval.

  % Example 3.46; the example itself uses calculus.
  \item For \(f(x)=x-x^3\) on \([-1,1]\), compute the image interval
    using calculus as in Example 3.46. Is it determined by the values
    at \(-1\) and \(1\)? What extra hypothesis would make
    \(f([a,b])=[f(a),f(b)]\) valid?

  % Examples 3.47--3.49.
  \item Give examples showing how the conclusion of Theorem 3.45
    can fail when continuity, the interval hypothesis, or compactness
    is omitted. In each case, compute the image and specify whether
    it fails to be an interval, to be closed, or to be bounded.

  % Final paragraph of Section 3.6.
  \item What can be said about the continuous image of an arbitrary
    interval? Must openness, closedness, or boundedness of the
    interval be preserved? Use examples from this chapter.
\end{enumerate}

\subsection{Monotonic functions}
\textit{Source: \S3.7, pp.~35--37.}

\begin{enumerate}[resume]
  % Definition 3.50.
  \item Define increasing, strictly increasing, decreasing, strictly
    decreasing, monotonic, and strictly monotonic functions on an interval,
    using the terminology in these notes.

  % Terminology discussion following Definition 3.50.
  \item Is a constant function increasing, decreasing, both, or neither
    under these conventions? Explain the alternative use of
    ``nondecreasing'' and ``nonincreasing'' and why conventions matter.

  % Theorem 3.51.
  \item State the one-sided-limit theorem for monotonic functions
    (Theorem 3.51). At which domain points does it guarantee finite
    one-sided limits?

  % Formulas in the proof of Theorem 3.51.
  \item For an increasing \(f:I\to\R\) and an interior point
    \(c\in I\), express its left and right limits as a supremum
    and an infimum of appropriate sets of values. How do these
    limits compare with \(f(c)\)?

  % Proof of Theorem 3.51 and its right-limit exercise.
  \item Prove the left-limit formula in Theorem 3.51 from completeness
    and monotonicity. Why is the relevant set of values nonempty
    and bounded above? Explain the corresponding right-limit argument
    and how to treat decreasing functions.

  % Endpoint discussion following Theorem 3.51.
  \item What one-sided limits are guaranteed for a monotonic function
    on a closed interval \([a,b]\), with \(a<b\)? Must they agree
    with the endpoint values? Give an example where they do not.

  % Corollary 3.52 and its proof.
  \item State and prove the restriction on interior discontinuities
    of monotonic functions (Corollary 3.52). Why can an interior
    discontinuity be neither removable nor essential?

  % Final paragraph of Section 3.7.
  \item What can be said about the number of discontinuities of a
    monotonic function? Must it be finite? Contrast this with the
    discontinuities of the Dirichlet function.
\end{enumerate}
